\section{Canonical transform and exact zero expansion}
\label{sec:explicit}

\subsection{The base-1 extension}

The uncut Mellin transform of \(dE\) is not defined in a right half-plane,
because \(dE=-dx\) on \((0,1)\).  Introduce only for the transform
\[
J_h(t)=
-\int_{[1,\infty)}
x^{-1/2}K_h(t-\log x)\,dE(x).
\]
If
\[
E_0(x)=\psi(x)-x,\qquad
E_1(x)=\psi(x)-(x-1)\quad(x\ge1),
\]
then \(E_0(1)=-1\), \(E_1(1)=0\), and distributionally
\[
d(\mathbf1_{[1,\infty)}E_0)
=\mathbf1_{[1,\infty)}dE-\delta_1,
\]
whereas
\[
d(\mathbf1_{[1,\infty)}E_1)
=\mathbf1_{[1,\infty)}dE.
\]
Thus \(J_h\) uses the endpoint-normalized primitive \(E_1\) and contains no
spurious \(K_h(t)\) boundary term.  It satisfies
\[
J_h(t)=0\quad(t\le-1),\qquad
J_h(t)=\Ih(t/2)\quad(t>1).
\]

\subsection{Bilateral Laplace transform}

For \(\Re s>1/2\), absolute Euler-product convergence and Tonelli's theorem
give
\begin{align}
\mathcal B_h(s)
&=
\int_\R J_h(t)e^{-st}\,dt
\notag\\
&=
\Lh(s)
\left\{
\frac{\zeta'}{\zeta}\left(s+\frac12\right)
+\frac{1}{s-\frac12}
\right\}.
\label{eq:B-transform}
\end{align}
At \(s=1/2\),
\[
\frac{\zeta'}{\zeta}\left(s+\frac12\right)
=-\frac{1}{s-1/2}+\gamma_E+O(s-1/2),
\]
so the apparent pole is removable and
\[
\mathcal B_h(1/2)=\gamma_E\Lh(1/2)=\gamma_E M_h^2.
\]
The remaining poles are
\[
s=\rho-\frac12
\]
for nontrivial zeros, with residue
\[
m_\rho\Lh\left(\rho-\frac12\right),
\]
and
\[
s=-2n-\frac12,\qquad n\ge1,
\]
for trivial zeros, with residue
\[
\Lh\left(-2n-\frac12\right).
\]
A multiple zero changes the residue by its multiplicity; it does not create
a derivative term.

\subsection{Contour displacement}

For \(c>1/2\), Fourier inversion on the line \(\Re s=c\) gives
\[
J_h(t)
=
\frac{1}{2\pi i}\int_{c-i\infty}^{c+i\infty}
e^{st}\Lh(s)
\left\{
\frac{\zeta'}{\zeta}\left(s+\frac12\right)
+\frac{1}{s-\frac12}
\right\}ds.
\]
The integral is absolutely convergent because, uniformly in every fixed
vertical strip,
\[
\Lh(\sigma+i\tau)=O_N((1+|\tau|)^{-N})
\]
for every \(N\).  Choose heights avoiding zero ordinates by
\(\gg1/\log T\); on those horizontal segments
\(\zeta'/\zeta=O(\log^2T)\), so the horizontal integrals vanish.

Move the left side to \(\Re s=-2N-1\).  Repeated integration by parts and
\(\supp K_h=[-1,1]\) give, for \(M\ge3\),
\[
|\Lh(-2N-1+i\tau)|
\le
C_Me^{2N+1}(1+N)^M(1+|\tau|)^{-M}.
\]
The remaining vertical integral is consequently
\[
O_{h,M}\left(
e^{-(2N+1)(t-1)}(1+N)^M\log(N+2)
\right),
\]
which tends to zero when \(t>1\).  Summing the crossed residues yields
\begin{align}
J_h(t)
={}&
\sum_\rho m_\rho
\Lh\left(\rho-\frac12\right)
e^{(\rho-1/2)t}
\notag\\
&+
\sum_{n=1}^{\infty}
\Lh\left(-2n-\frac12\right)e^{(-2n-1/2)t}.
\label{eq:J-explicit}
\end{align}
Since \(\Lh\) is even, \(t=2y\) proves
\eqref{eq:I-expansion}.

The trivial-zero sum is absolutely convergent for \(t>1\).  The zero-counting
bound \(N_\zeta(T)=O(T\log T)\), combined with rapid vertical decay of
\(\Lh\), proves absolute and locally uniform convergence of the nontrivial
zero sum after any fixed number of real derivatives.

\subsection{Archimedean cancellation}

For the reflected cross correlation, \(C(0)=0\), \(C(r)=0\) for \(r\ge0\),
and \(C(-r)=K_h(2y-r)\).  Equation~\eqref{eq:real-place} gives
\[
A_h^\infty(y)
=-\int_0^\infty
\frac{e^{r/2}}{e^r-e^{-r}}K_h(2y-r)\,dr.
\]
Using
\[
\frac{e^{r/2}}{e^r-e^{-r}}
=\sum_{n=0}^{\infty}e^{-(2n+1/2)r},
\]
Tonelli's theorem and \(u=2y-r\) give
\[
A_h^\infty(y)
=-\sum_{n=0}^{\infty}
\Lh\left(2n+\frac12\right)e^{-(4n+1)y}.
\]
The \(n=0\) term is \(-M_h^2e^{-y}\); the terms \(n\ge1\) cancel the
complete trivial-zero series in \eqref{eq:I-expansion}.  Combining with
\eqref{eq:Q-decomposition} proves the zero-only formulas
\eqref{eq:zero-only} and \eqref{eq:delta-zero-only}.
